
Perkembangan teknologi informasi saat ini menuntut adanya integrasi sistem digital di lingkungan institusi pendidikan guna meningkatkan efisiensi administrasi akademik. Namun, pelaksanaan Praktik Kerja Lapangan (magang) di SMK PGRI Telagasari masih dihadapkan pada kendala manajerial yang signifikan akibat ketergantungan pada prosedur konvensional. Proses pencatatan kehadiran, pengisian aktivitas harian siswa (logbook), serta rekapitulasi nilai masih bergantung penuh pada dokumen fisik. Keterbatasan ini menimbulkan berbagai risiko operasional seperti keterlambatan pelaporan, risiko kehilangan data, lambatnya verifikasi administratif, yang menyulitkan pihak manajemen sekolah dalam melakukan evaluasi cepat yang krusial bagi pengambilan keputusan strategis.

Sebagai solusi atas permasalahan tersebut, proyek akhir ini mengembangkan sebuah sistem informasi manajemen monitoring PKL bernama InternSync. Sistem ini dibangun menggunakan arsitektur client-server berbasis REST API terpusat dengan mengintegrasikan teknologi modern, meliputi React pada sisi frontend web, framework Laravel pada sisi backend, dan basis data MySQL. Pengembangan aplikasi seluler berbasis Flutter masih berada pada tahap perencanaan sebagai pengembangan lanjutan. InternSync menyediakan manajemen hak akses pengguna (Role-Based Access Control) untuk lima peran utama sekolah serta akses token aman bagi pihak industri. Sistem ini mengintegrasikan katalog lowongan PKL yang dikelola Koordinator dan dipilih oleh siswa dengan seleksi oleh Koordinator, jurnal aktivitas harian (logbook) digital, sistem perizinan presensi dengan konfirmasi industri fleksibel berfasilitas fallback sekolah, serta evaluasi akhir berbasis rubrik baku. Pengujian sistem InternSync dilakukan secara komprehensif menggunakan metode Black Box Testing yang berfokus pada validasi fungsionalitas input dan output pada lingkungan deployment publik. Hasil pengujian menunjukkan bahwa seluruh parameter fungsionalitas inti, mencakup operasi CRUD pada modul pengajuan, manajemen kuota mitra industri, pemantauan logbook, mekanisme penanganan kesalahan (error handling), serta responsivitas antarmuka web pada berbagai ukuran layar, berjalan sukses. Sistem InternSync terbukti andal dalam mempercepat alur birokrasi dan mengoptimalkan fungsi pengawasan monitoring PKL secara real-time demi mendukung implementasi konsep Smart School di SMK PGRI Telagasari.

\textbf{Kata Kunci:} Sistem Informasi Manajemen, Monitoring PKL, Laravel, React, REST API.
